\documentclass[11pt]{article}
\usepackage[margin=1in]{geometry}
\usepackage{amsmath,amssymb}
\usepackage{hyperref}
\usepackage{enumitem}
\usepackage{listings}
\usepackage{xcolor}
\usepackage{graphicx}
\usepackage{setspace}
\setstretch{1.15}

\definecolor{codegray}{gray}{0.95}
\lstset{
    backgroundcolor=\color{codegray},
    basicstyle=\ttfamily\small,
    breaklines=true,
    frame=single
}

\title{Brownian-Dynamics Kern--Frenkel EPYC1--Rubisco GPU Model: Changes and Detailed Description}
\author{}
\date{}

\begin{document}
\maketitle
\tableofcontents

\section{Purpose of the Rewrite}
The previous GPU rewrite captured the scaffold--linker architecture and Kern--Frenkel-style geometric gating, but it did not produce sufficiently dynamic-looking assemblies in trajectory space. Clusters could grow in terms of connectivity while appearing spatially static or insufficiently fluid in OVITO. The main reason was that the earlier implementation emphasized bond bookkeeping more strongly than force-driven structural relaxation.

The Brownian-dynamics rewrite was introduced to address exactly this limitation while preserving the main scientific ingredients:
\begin{itemize}
    \item explicit Rubisco scaffold particles,
    \item explicit EPYC1 chains,
    \item sticker-mediated binding,
    \item explicit Rubisco patches,
    \item radial and angular Kern--Frenkel gating,
    \item one-sticker/one-patch occupancy,
    \item reversible binding,
    \item timecourse outputs and OVITO trajectories.
\end{itemize}

\section{What Changed}
Relative to the prior KF GPU model, the Brownian rewrite introduces:
\begin{enumerate}[label=\arabic*.]
    \item chain spring forces,
    \item bound sticker--patch spring forces,
    \item short-range repulsions,
    \item force-based translational motion,
    \item gradual rotational diffusion.
\end{enumerate}

It preserves:
\begin{enumerate}[label=\arabic*.]
    \item patch-based Rubisco geometry,
    \item EPYC1 chains,
    \item Kern--Frenkel radial shell,
    \item Kern--Frenkel angular gating,
    \item occupancy constraints,
    \item reversible kinetics,
    \item CSV/plot/OVITO outputs,
    \item the same sweep regions.
\end{enumerate}

\section{System Components}
\subsection{Rubisco scaffolds}
Rubisco is represented as rigid scaffold particles with center positions $\mathbf{R}_i$, orientation matrices $\mathbf{Q}_i$, and explicit patches. Patch positions are obtained by rotating reference patch coordinates:
\[
\mathbf{p}_{i\alpha} = \mathbf{R}_i + \mathbf{Q}_i \mathbf{p}^{(0)}_{\alpha}.
\]

\subsection{EPYC1 chains}
Each EPYC1 molecule is a bead-spring chain with pattern
\[
S-T-S-T-S-T-S-T-S.
\]
Only the $T$ beads are stickers capable of binding Rubisco patches.

\section{Kern--Frenkel Criterion}
A sticker can bind a patch only when two conditions are met.

\subsection{Radial shell}
\[
\sigma < r < \sigma \lambda.
\]

\subsection{Angular condition}
If $\hat{\mathbf{n}}_p$ is the outward patch normal and $\hat{\mathbf{d}}$ is the unit vector from patch to sticker, then:
\[
\hat{\mathbf{n}}_p \cdot (-\hat{\mathbf{d}}) > \cos\theta_c.
\]

This enforces directional specificity.

\section{Brownian Dynamics}
The key new feature is that positions evolve under both deterministic forces and stochastic noise:
\[
\mathbf{x}(t+\Delta t)=\mathbf{x}(t)+\mu \mathbf{F}(t)+\boldsymbol{\eta}.
\]

\subsection{Chain spring forces}
Adjacent EPYC1 beads are connected by harmonic springs:
\[
\mathbf{F}_{ij}^{\text{chain}} = -k_{\text{chain}}(r_{ij}-r_0)\hat{\mathbf{r}}_{ij}.
\]

\subsection{Bound sticker--patch forces}
A bound sticker and occupied patch are coupled by an effective spring:
\[
\mathbf{F}_{\text{bind}} = -k_{\text{bind}}(r-r_{\text{bind},0})\hat{\mathbf{r}}.
\]

\subsection{Repulsion}
Rubisco--Rubisco and bead--Rubisco repulsions prevent unrealistic overlap.

\subsection{Rotational diffusion}
Rubisco orientations are updated using small random rotations:
\[
\mathbf{Q}_i(t+\Delta t)=\Delta\mathbf{R}_i\,\mathbf{Q}_i(t).
\]

\section{Outputs}
For each parameter set the code writes:
\begin{itemize}
    \item CSV timecourses with
    \begin{itemize}
        \item step,
        \item time,
        \item sigma,
        \item lambda,
        \item eta,
        \item number of clusters,
        \item largest cluster size,
        \item largest cluster fraction,
        \item number of bonds.
    \end{itemize}
    \item JPEG plots of each of those main observables.
    \item XYZ trajectories with particle types:
    \begin{itemize}
        \item R = Rubisco,
        \item S = spacer,
        \item T = sticker.
    \end{itemize}
\end{itemize}

\section{Why It Should Look More Dynamic}
The previous KF version mainly changed connectivity. The Brownian version adds force-driven structural relaxation. Therefore clusters should now:
\begin{itemize}
    \item move more fluidly,
    \item compact after bond formation,
    \item reorganize under linker-mediated forces,
    \item diffuse under stochastic noise while maintaining structure.
\end{itemize}

\section{Interpretation}
The purpose of the rewrite is to move the model from
\begin{quote}
correct connectivity growth
\end{quote}
to
\begin{quote}
correct connectivity growth plus physically plausible spatial dynamics.
\end{quote}

That is the main conceptual difference from the earlier GPU KF implementation.

\end{document}
